\section{Major Risks / Reality Check}
\label{sec:risks}

Category legend used in the table: (1) fundamentally impossible, (2) technically difficult, (3) commercially difficult, (4) solvable with customer engineering data, (5) solvable through Alfred hardware, (6) solvable through Alfred software. Most real risks are more than one category at once; the table lists the \emph{dominant} category and explains the rest in prose.

\begin{center}
\small
\begin{longtable}{@{}p{2.4cm} p{1.7cm} p{11.2cm}@{}}
\toprule
\textbf{Limitation} & \textbf{Category} & \textbf{Assessment} \\
\midrule
\endhead
Encrypted traffic (SecOC/MACsec-protected signals) & 2, 4 & Not fundamentally impossible in principle (Alfred is not attempting to defeat the crypto), but a mapping cannot be validated through Structured Experiment Mode if the payload is encrypted and Alfred has no key. Solvable only with customer-provided keys/engineering data under authorization, or by treating the encrypted signal as opaque and mapping around it (e.g., inferring intent from an unencrypted companion signal). Do not budget on Alfred cracking this independently. \\
Authenticated diagnostics (UDS security access) & 2, 4 & Read-only diagnostic services are frequently accessible without security access; programming/actuation-class services are gated by a seed-key exchange Alfred will not attempt to derive (Section~\ref{sec:deployment-reprogramming}). Solvable when the OEM/Tier-1 provides or authorizes the algorithm; otherwise those specific services are simply out of reach, which is an acceptable, bounded limitation, not a blocker to the broader architecture. \\
Proprietary protocols (ag/mining/industrial, no public spec) & 2, 6 & This is exactly what Structured Experiment Mode and the Correlation Engine are for (Section~\ref{sec:probe-discovery}); technically difficult and slower than DBC-assisted discovery, but not blocked --- it is Alfred's primary differentiator in these verticals precisely because it is hard. \\
Gateways (filtering/segmenting traffic) & 2, 4, 5 & A single-point probe cannot see traffic a gateway filters; distributed probe placement on both sides of a gateway (Section~\ref{sec:probe-hardware}) is the hardware mitigation, and customer-provided gateway routing tables (engineering data) resolve ambiguity faster than blind multi-point capture alone. \\
Switched automotive Ethernet visibility & 2, 5 & A switch does not broadcast unicast traffic to arbitrary ports; passive visibility requires either a true inline tap (Alfred T1 TAP, Section~\ref{sec:product-family}) at each segment of interest, or a SPAN/mirror-capable switch port if the vehicle's own switch supports it (rarely reliable to assume). This is a hardware placement problem, solved by deploying enough taps, not a software problem. \\
Inaccessible local SPI/I2C buses (inside a sealed ECU) & 1 (for that ECU, non-invasively), 5 & If the bus never leaves the ECU's own PCB, it is not observable without physically opening/instrumenting that specific ECU, which is a bench/lab (not in-vehicle) activity requiring the physical unit and is out of scope for standard in-vehicle probing. Genuinely unreachable for a non-invasive engagement; only solvable by Alfred hardware if the customer supplies the ECU for bench-level reverse engineering as a distinct, explicit workstream. \\
ECU secure boot / cryptographic root of trust & 1 (bypass), 5/6 (adapter path) & Bypassing this is out of scope categorically (Section~\ref{sec:deployment-reprogramming}). The answer is never "defeat it," it is the Legacy Adapter/Proxy pattern: leave the ECU's firmware untouched and command it through its existing legitimate interface instead. \\
Safety requirements for actuated functions & 2, 3 & Genuinely hard, not a paperwork problem: closing the loop on a real safety case (ISO~26262 ASIL classification, FMEDA, safety manual) for any function Alfred actuates is significant engineering and, for anything beyond QM/ASIL~A, likely requires a dedicated functional safety engineer and independent assessment --- budget real headcount and calendar time here, not just "add tests." \\
Real-time control latency over Ethernet & 2, 6 & Bounded and solvable with correct architecture (TSN traffic shaping, correct placement per Section~\ref{sec:acc}'s framework, keeping hard-real-time loops local) but genuinely requires getting the placement decision right per control loop; a placement mistake (centralizing something that needed edge execution) manifests as a real, hard-to-retrofit latency bug, not a tunable parameter. \\
Network bandwidth (backbone contention) & 2, 6 & 10BASE-T1S's shared-medium PLCA arbitration caps aggregate spur throughput; this is a capacity-planning problem solved by correct network segmentation (Section~\ref{sec:option-a-overview}) and TSN prioritization, not by faster hardware alone --- a design that puts too many nodes on one T1S spur will hit a real ceiling. \\
Determinism (jitter under real OS/software load) & 2, 6 & Achievable on the RT domain (Section~\ref{sec:acc}) but requires actual RTOS scheduling discipline (bounded interrupt latency, WCET analysis) --- this is standard hard-real-time embedded engineering rigor, not a novel problem, but it is easy to erode accidentally by adding "just one more" software feature to the RT domain over time. \\
Clock synchronization accuracy & 2, 5, 6 & gPTP is a mature, solved protocol, but achieving it reliably requires hardware timestamping support in every hop (Section~\ref{sec:probe-discovery}) --- a cheap non-timestamping switch/NIC in the chain silently degrades sync quality; this is a component-selection discipline problem, not a protocol problem. \\
Electrical compatibility (probe/E2B on live automotive buses) & 2, 5 & Getting transceiver front ends, protection, and isolation right (Section~\ref{sec:bom}) is well-understood engineering, but under-protecting a probe front end risks disturbing the very bus being observed (or damaging the probe) --- treat probe electrical design with the same rigor as production hardware even though the unit itself is "just" a dev tool. \\
Automotive EMC (radiated/conducted emissions and immunity) & 2, 3 & Full EMC qualification (CISPR 25, ISO 11452) is expensive and slow; acceptable to defer for bench-only/development-vehicle prototypes (MVP0--2) but is a real, multi-month, real-cost gate before any production-installed hardware (Gateway, E2B) ships in volume. \\
ISO~26262 (functional safety) & 2, 3 & Not a checkbox; requires process discipline (Part 6/8 work products) from early in development for anything safety-relevant, which is exactly why Section~\ref{sec:firmware-generation}'s pipeline is built to emit certification evidence as a byproduct rather than as an afterthought --- retrofitting evidence generation after the fact is materially more expensive than building it in. \\
ISO/SAE~21434 (cybersecurity engineering) & 2, 3, 6 & Threat analysis and risk assessment (TARA) obligations apply to anything Alfred connects to a vehicle network, including the Probe/Gateway itself; Alfred's own hardware becomes an attack-surface item in the \emph{customer's} TARA, which must be addressed with real security engineering (secure boot, signed firmware, the transmit-interlock hardware in Section~\ref{sec:active-control}), not assumed away because "it's just a probe." \\
ASPICE (process maturity assessments) & 3 & Primarily a commercial/procurement gate for OEM/Tier-1 engagements, not a technical blocker; large customers may require ASPICE capability-level evidence before a production contract regardless of the technology's merit --- budget this as a sales-cycle cost, addressed with process tooling once a program reaches that stage, not before. \\
Silicon qualification (AEC-Q100/Q200) & 2, 3, 5 & Straightforward in principle (choose qualified parts, Section~\ref{sec:bom}) but real in cost and lead time --- automotive-grade parts have longer lead times and higher unit cost than the prototype equivalents, and a BOM designed around prototype parts is not a drop-in swap later without a board respin. \\
Hardware cost at low volume & 3 & E2B/Probe unit cost at prototype/pilot volumes will be materially higher per unit than at scale; this is a normal hardware-startup economics problem, not specific to Alfred's architecture, but should be modeled explicitly before quoting pilot pricing. \\
Harness topology (retrofit fit) & 2, 3, 5 & Every vehicle's physical harness routing and connector placement is different; an E2B node replacing a legacy ECU (Phase 3 migration, Section~\ref{sec:migration}) must physically fit and connect where the old part was, which constrains E2B mechanical/connector design per program and is a real integration engineering cost each time, not a one-time solved problem. \\
Connector complexity & 2, 5 & Standardizing on one connector family across the E2B line (Section~\ref{sec:e2b-family}) helps but does not eliminate the need to support whatever connector the \emph{peripheral} side already uses in a retrofit scenario --- expect a real, ongoing library of peripheral-side connector/pigtail adapters as a cost of doing brownfield retrofits. \\
\bottomrule
\end{longtable}
\end{center}
